%==============================================================================
\section{Protocol: claims committed before code}
\label{sec:protocol}

The kernel was built in stages, and each stage after the first followed one protocol.
First, the claims the stage would test were written in a plain-text record, each with a statement of the outcome that would count against it.
Where a control was needed to make the test able to fail, the record named it too.
Second, the record was committed to version control.
Third, the code that tested the claims was written and run, and its outcome was recorded in a separate file that never edits the first.
The commit history therefore orders each claim before its test \citep{nosek2018preregistration}.

Three kinds of entry are distinguished in the record.
\emph{Predictions} state a direction and a threshold in advance.
\emph{Measurements} are instruments whose numbers are reported but not predicted, such as the census of gate letters; they are not counted as claims.
\emph{Reproductions} re-run an analysis that had been run before the record was written; they are labelled as such and are not counted either.

The record uses short identifiers that repeat across stages.
\Cref{tab:ids} maps them to the identifiers used in this paper.
Each claim's full text is reproduced in \cref{app:claims}.

\begin{table}[t]
\centering
\caption{Claims tested, with the commit that recorded each before its code was written.
``Record'' gives the identifier in the repository's prediction file.
Each stage name refers to a section of the paper.}
\label{tab:ids}
\footnotesize
\setlength{\tabcolsep}{4pt}
\begin{tabular}{@{}lp{3.6cm}p{3.4cm}ll@{}}
\toprule
Stage & This paper & Record & Recorded in & Kind \\
\midrule
Medium (\cref{sec:res-medium}) & M1; M2a, M2b; M3a, M3b & M1; M2; M3 & \texttt{de6f210} & predictions \\
Kinetics (\cref{sec:res-kinetics}) & K1 & P1 & \texttt{de6f210} & prediction \\
Folding (\cref{sec:res-folding}) & F1 & F1 & \texttt{d77d09c} & reproduction \\
 & F2; F3; F4a, F4b & F2; F3; F4 & \texttt{d77d09c} & predictions \\
 & F5 & F5 & \texttt{d77d09c} & measurement \\
Extensions (\cref{sec:res-extend}) & E1--E3 & X1--X3 & \texttt{cddf153} & measurements \\
 & E4a, E4b, E4c; E5; E6 & X4a, X4b, X4c; X5; X6 & \texttt{cddf153} & predictions \\
Databases (\cref{sec:res-db}) & R1, I1, L1, S1, C1, A1, D1, G1 & R1, X1, L1, S1, C1, P1, D1, G1 & \texttt{b1d7fc8} & predictions \\
\bottomrule
\end{tabular}
\end{table}

\paragraph{Statistics.}
Paired per-reaction comparisons use the Wilcoxon signed-rank test \citep{wilcoxon1945individual}; comparisons of two samples use the Mann--Whitney test \citep{mann1947test}, one-sided where the claim states a direction.
Agreement between partitions uses the adjusted Rand index \citep{hubert1985comparing}.
Label and signature tests use empirical $p$-values from 200 permutations of the labels.
The threshold is $\alpha=0.05$ throughout, as written in the record.
No correction for multiple testing was declared in advance, and none is applied; \cref{sec:limitations} discusses the consequence.

\paragraph{Counting outcomes.}
A claim with two parts that can fail separately is counted as two (M2, M3, F4, E4).
S1 also has two parts, but its record states a single verdict for both, so it is counted once.
E4c could not be tested because the kinetic constant it needs is not in SABIO-RK; it is counted as untestable, not as a failure, as its record specifies.
The resulting 23 claims are listed with their outcomes in \cref{tab:outcomes}.
